
Assume the default: things are relatively boring. They are stuck together in a fixed pattern of interaction, like asex (or maybe even like rocks). Then, there's some innovation that randomly (or intentionally, later in evolution) places a barrier to limit fixed association. The larger system (including that barrier) becomes more generative, and therefore more successful. 

Is it always the case that there's an `outer' optimization pressure (like group selection) and an `inner' optimization pressure (like individual selection)? Does this mean that there are also always `outer' and `inner' barriers? Like an outer barrier that separates the two optimization pressures, and an inner barrier that results from the outer pressure? 

So in sex, we have the 

[...and, how does this work with awareness? ah, there's the local perspective (inner optimization), and whatever is outside of it (outer optimization). there's something about how this bootstraps itself into a fractal.]

With sex, the optimization pressure is at the level of meta-selection (not selecting for an organism but for a population). All the genes are stuck together in a fixed association, and base selection is selecting for that. 

How does this work with drives? If you have an unchecked drive (like in overeating), it's a detrimental fixed pattern. The optimization pressure is at the level of organism-evolution. A single drive is a proxy for the organism's fitness. Evolution finds mechanisms to limit the effect of that drive (satiety, competing drives etc). However, there's also a very different mechanism where 

How about strategies? Hershberger's chickens had an unchecked strategy (another inflexible pattern). But now the optimization pressure is at the level of the drive. If the strategy doesn't work (the drive is stymied through that route), then the *pressure of the drive*  We have mechanisms like arbitration between decision systems to 
But is it the world, or an internal barrier, that stymies the drive? Ultimately, it's an internal barrier. Because Hershberger's chickens didn't have that barrier, and they went on pursuing the strategy. 

How about groups? This is where humans invent a strategy as a barrier. 

In ecosystems, it's back to something like group selection. More similar to sex. 


overfitting perspective. a gene becomes overfit to its context. through some optimization pressure. 

Genomes.
Core idea: there's a collapse to 
Or, with genes, it's not really a collapse because it's not like a richer system was built up and it risks collapsing....
Collapsed configuration (easy way out): selfing/asexual. fixed genome. each gene becomes overfit to particular genetic context. 
The drive to reproduce is "trying" to do something.
Or is it the particular genetic configuration that's "trying" to do something?

Ecosystems.
Core idea: there's a collapse to fewer species.

Drives and goals.
Core idea: there's a collapse to a single drive. 
The drive to eat is `trying' to do something. It has a myopic goal. If that goal is allowed to 

Frames.
Core idea: there's a collapse to a particular way of thinking.
Take old-school religion, for example. Believing in a particular God and that worshippers of other gods are evil. It has a myopic goal. If left unchecked, it leads to erasure of diversity. 

Groups. 
Core idea: there's a collapse to groupthink.
Now we have ideas in the group. It's less clear that the *ideas* per se have myopic goals. An individual in the group has a wish to solve the problem. If short-circuited (not having to struggle against the problem), they might find a 

%----------------------------------------------------------------------
% Concept x example matrix (Table 1). Rows = all examples, including the
% minor ones; last row bridges to Section 2. Cells marked with a dagger
% are proposals that go beyond what the paper/outline currently say.
%----------------------------------------------------------------------

\begin{landscape}
\begin{table}[p]
\centering
\tiny
\setlength{\tabcolsep}{3pt}
\renewcommand{\arraystretch}{1.15}
\caption{\textbf{Concept $\times$ example matrix, Part I: main examples.} Rows are the paper's worked examples (Part II continues with the minor examples from \S1 ``Other examples'' and a row bridging to \S2). The columns form one causal arc, read left to right: a stored potential discharges along its congruent path (\textit{short-circuiting}) into a collapsed state (\textit{fixation}), unless a \textit{semi-permeable boundary} blocks that path; the boundary \textit{contextualizes} the pattern rather than nullifying it, which gives the parts \textit{modularity}, gives the adaptive process \textit{evolvability}, and gives the whole open-ended \textit{discovery}. Cells marked $^{\dagger}$ are proposals that go beyond the current draft and outline.}
\label{tab:conceptmap}
\medskip
\begin{tabularx}{\linewidth}{@{}>{\RaggedRight\arraybackslash\bfseries}p{1.9cm} Y Y Y Y Y Y Y@{}}
\toprule
 & \textbf{Fixation} (collapsed state) & \textbf{Short-circuiting} (discharge dynamic) & \textbf{Semi-permeable boundary} (mechanism) & \textbf{Contextualization} (the boundary's action) & \textbf{Modularity} (what the parts gain) & \textbf{Evolvability} (what adaptation gains) & \textbf{Discovery} (what the whole gains) \\
\midrule
Genomes &
Alleles overfit to one fixed genetic background; deleterious mutations permanent (Muller's ratchet); beneficial mutations in separate lineages can never combine. &
$^{\dagger}$Selfing/asexuality is the congruent path to replication: cheap copies now, while stored variation (heterozygosity) literally decays; meiotic drivers cheat fair segregation for the same trade. &
Recombination, self-incompatibility, inbreeding avoidance --- barriers placed directly on the path to reproduction. &
A gene's selfishness is held in check by the whole genome; each allele is expressed as one contribution among shifting backgrounds. &
Alleles selected across many backgrounds become interchangeable units that contribute to any context (mixability). &
Selection acts on loci semi-independently; beneficial variants become combinable; adaptation accelerates (evolvability as meta-learning). &
Novel genotypes assembled from parts proven in separate lineages; introgression yields some of the fastest evolution observed. \\
\addlinespace
Drives \& goals &
One drive or prepotent strategy dominates: obesity, addiction, burnout; Hershberger's chickens beaten by ``approach food.'' &
The drive discharges along its congruent path, consuming the proxy rather than the fitness it stands for (wireheading in the limit). &
Cognitive control: context, rules and distal goals bias the competition among drives and responses. &
Expression made conditional, not nullified: hunger guides eating without overeating; absorption in work coexists with a bedtime rule. &
$^{\dagger}$Drives and strategies become a repertoire of situational modules rather than monolithic controllers of behavior. &
$^{\dagger}$Strategies can be revised without unlearning the underlying drive; conditional modules recombine into novel plans. &
Blocking the congruent path forces symmetry breaking; stymying a shallow objective yields progress on a deeper one (insight, creativity). \\
\addlinespace
Frames &
Rigid negative beliefs, obsession, delusion; scientific dogma. From inside, the frame is invisible and mistaken for reality. &
$^{\dagger}$The frame confirms itself: evidence is assimilated, anomalies are discharged as noise, and inquiry never has to leave the frame. &
Awareness and metacognition; in science, anomaly-tracking, peer critique, calibrated resistance to change (Kuhn, Popper). &
The frame is retained as a partial, evidence-responsive model: committed to provisionally, not absolutely. &
$^{\dagger}$Theories become portable instruments applicable across domains; several frames can be held and deliberately switched. &
Anomaly $\to$ crisis $\to$ revolution: the community can restructure its own descriptions of the world. &
Paradigm shifts; reframing as the core mechanism of insight. \\
\addlinespace
Groups &
Homogenization: information cascades, groupthink; the peptic-ulcer consensus wrong for decades. &
$^{\dagger}$Continuous sharing lets seductive early solutions propagate instantly; copying discharges the group's distributed search before it runs. &
Intermittent communication; independent work then aggregation (Craven); source tagging; protected dissent; adversarial scrutiny. &
Each view enters as one signal among many (e.g., Bayesian aggregation), influencing the group without overwriting anyone. &
Members develop distinct, complementary expertise --- composable components rather than copies of one another. &
$^{\dagger}$The group maintains a population of diverse solutions, so it can adapt when the problem shifts. &
Compositional recombination: steel punch $+$ screw press $\to$ printing press. \\
\addlinespace
Institutions &
One actor's agenda captures the system: monopoly, domination, regulatory capture. &
$^{\dagger}$The motive gratified by the direct route: profit via fraud or suppression of competition, power via silencing --- extraction instead of creation. &
Constitutions, laws, norms, courts, regulators. &
Motives rerouted, not annulled: profit under competition and regulation fuels innovation. &
$^{\dagger}$Separation of powers; firms, agencies and jurisdictions as bounded, specialized units (polycentric governance). &
Institutions adapt as loopholes are found, gaining robustness from being tested against many situations and motives. &
Rerouted energy generates new products, markets and institutional forms. \\
\bottomrule
\end{tabularx}
\end{table}

\begin{table}[p]
\centering
\tiny
\setlength{\tabcolsep}{3pt}
\renewcommand{\arraystretch}{1.15}
\caption{\textbf{Concept $\times$ example matrix, Part II: minor examples and the bridge to \S2.} Columns as in Part I.}
\label{tab:conceptmap2}
\medskip
\begin{tabularx}{\linewidth}{@{}>{\RaggedRight\arraybackslash\bfseries}p{1.9cm} Y Y Y Y Y Y Y@{}}
\toprule
 & \textbf{Fixation} (collapsed state) & \textbf{Short-circuiting} (discharge dynamic) & \textbf{Semi-permeable boundary} (mechanism) & \textbf{Contextualization} (the boundary's action) & \textbf{Modularity} (what the parts gain) & \textbf{Evolvability} (what adaptation gains) & \textbf{Discovery} (what the whole gains) \\
\midrule
Cell membranes &
Equilibrium: gradients collapsed, the cell indistinguishable from its surroundings --- death as literal fixation. &
The literal case: ions discharging straight down their gradient through an unguarded membrane, annihilating the cell. &
The membrane itself, with conditional gates, channels and pumps. &
The gradient that would destroy the cell is put to work: ATP synthesis, action potentials. &
Compartments (cells, organelles), each maintaining a distinct internal chemistry. &
$^{\dagger}$New gate types add capabilities without rebuilding the cell; compartments vary independently. &
Gated gradients underlie signaling and nervous systems --- new function from harnessed potential. \\
\addlinespace
Symbiogenesis &
Each lineage confined to its own fitness well for hundreds of millions of years. &
$^{\dagger}$Engulfment's congruent path is digestion: consuming the partner as food and dissipating the encounter's potential. &
Long lineage isolation beforehand; afterward, retained membranes keeping the endosymbiont distinct. &
The bacterium keeps its own genome and membrane yet functions as an organelle within the host's economy. &
The mitochondrion as a bolt-on power module; the two genomes remain semi-distinct. &
Solutions found in isolation become combinable, opening the eukaryotic radiation. &
A fundamentally new kind of entity that neither lineage could reach alone. \\
\addlinespace
Brains &
Hypersynchrony: seizure; pathological correlation in Parkinson's; blurred representations that interfere. &
$^{\dagger}$Runaway excitation recruiting the whole network --- activity spreading down the unguarded path (the seizure as literal short circuit). &
Lateral inhibition, pattern separation, excitatory--inhibitory balance, oscillatory gating. &
Representations kept distinct yet interacting under control, between pure synchrony and uncorrelated noise. &
Pattern-separated memories and specialized circuits as discrete, reusable units. &
New learning without catastrophic forgetting; separated elements remain recombinable. &
Replay and imagination: structured synthesis composing sequences never experienced. \\
\addlinespace
Relationships &
Enmeshment: whose experience is whose blurs; one perspective absorbs the other. (Rigid walls are the opposite collapse: isolation.) &
$^{\dagger}$Discomfort discharged into the other --- reassurance-seeking, merging, conforming --- instead of being metabolized into one's own structure. &
Psychological boundaries: conditional openness about where self ends and other begins. &
The other's experience is received as theirs: informative without being commandeering. &
Each person remains a distinct self, able to participate in many relationships and roles. &
$^{\dagger}$Partners keep developing rather than freezing into complementary roles; the relationship can renegotiate itself. &
Genuine dialogue produces understanding neither person held alone. \\
\addlinespace
Ecosystems &
Monoculture; an invasive species overwhelming the community; trophic collapse. &
$^{\dagger}$A population consuming its resource base unchecked: boom and crash; the invader with no predators takes the congruent path to collective collapse. &
Predation, competition, parasitism holding every population in check. &
Each species' expansion is held to a niche, where its consumption becomes a structuring role (Yellowstone wolves). &
Niches and functional guilds as composable units of community structure. &
Diversity as a reservoir; coevolution continually opens new adaptive paths. &
New niches and functions continually created (ecosystem engineers: beavers, reefs). \\
\addlinespace
Machine learning (bridge to \S2) &
Overfitting; mode collapse; co-adapted units; policy collapse onto a proxy objective. &
Reward hacking: optimizing the measurable proxy directly (maximize `likes' $\to$ flattery); shortcut features (elephant $\to$ giraffe). &
Dropout, regularization, KL penalties, withheld test sets, stop-gradients, limited inter-agent communication, sandboxing. &
$^{\dagger}$The proxy objective still steers learning but cannot be perfectly satisfied, forcing solutions that track deeper structure. &
Ensembles, mixtures of experts, slot attention: components forced to be separately competent. &
Meta-learning; population-based training preserving diversity for continued adaptation. &
Generalization beyond the training distribution; novelty search; capabilities not explicitly specified. \\
\bottomrule
\end{tabularx}
\end{table}
\end{landscape}

\note{Unification the table suggests, re the open fixation-vs-short-circuiting questions in \S2.1: every row contains a stored potential (an ion gradient, a drive, heterozygosity, a group's distributed search, an actor's motive, a partner's inner life). \textit{Short-circuiting} is the \textbf{event}: the potential discharging along its congruent, lowest-resistance path. \textit{Fixation} is the \textbf{state}: the collapsed arrangement left behind --- or the standing lock-up when structure freezes without any discharge at all (linkage). On this reading the two are not rival framings; short-circuiting is the dynamics whose endpoint is fixation. For non-homeostatic artifacts (the train brakes), where nothing is stored, fixation/overfitting is the general concept and short-circuiting is the special case that appears whenever potential is stored --- matching the hunch in the draft. The membrane row makes the metaphor literal and could anchor the terminology.}

\note{Second observation from filling in the matrix: the seven columns are not seven independent concepts but three groups --- \textbf{failure} (fixation as state, short-circuiting as dynamics), \textbf{mechanism} (semi-permeable boundary, whose action is contextualization), and \textbf{payoff at three scales} (part: modularity; adaptive process: evolvability; whole: discovery). If space is tight, the table could be collapsed to failure / mechanism / payoff.}